
    \section{Conclusiones}
En este trabajo se montó y caracterizó un reactor de plasma para la remediación de aguas contaminadas, optimizando su diseño y eficiencia en la eliminación del azul de metileno. En particular, se mejoró el sistema de electrodos con respecto a trabajos anteriores, logrando un diseño que minimiza la ocurrencia de descargas secundarias no deseadas. Además, se perfeccionó el método de medición de potencia eléctrica y se caracterizó el comportamiento de distintos dieléctricos como acrílico, teflón y vidrio. En las experiencias realizadas, se llegó a degradar el $(90 \pm 2) \%$ del azul de metileno tras un tratamiento de 140 minutos con un rendimiento energético de $(1,0 \pm 0,5) \n{g}/\tx{kWh}$. Por otro lado, los ensayos combinando el tratamiento con plasma y recubrimientos de $\tx{TiO}_2$ (reactor híbrido) presentaron una eficiencia de degradación menor en comparación a los obtenidos en el reactor de plasma. Para futuros trabajos se propone agregar al reactor una lámpara externa de radiación ultravioleta que estimule la acción del dióxido de titanio, estudiar distintos materiales sobre los cuales realizar el recubrimiento (como acero inoxidable), o disminuir su grosor para lograr una mayor eficiencia de degradación.
